\begin{table}
\centering\small
\begin{tabularx}{\linewidth}{lp{2.9cm}S[table-format=3.2]S[table-format=3.2]Xp{2.3cm}}
\toprule
Ítem & Partida & {Antes} & {Ahora} & Motivo & Detectado por\\
\midrule
2.4 & Relleno estabilizado e = 0.15 m & \num{9.36} & \num{7.65} & Relleno sobre la huella; EETT 2.4 lo define como base del radier (interior) & Revisor\\
11.1 & Cerámica 30 x 30 & \num{51.03} & \num{7.95} & Cerámica en todo el interior; EETT 11.2: solo cocina y baño & Revisor\\
3.3 & Tensores & \num{164.84} & \num{136.27} & Doble conteo de $\phi$8; faltaban +5\,\% de patas y la penetración de 40\,$d$ en fundación (NCh 353 6.2.8) & Revisor, ciega, NCh 353\\
3.2 & Escalerilla & \num{100.16} & \num{67.12} & Sin descontar vanos y con alambre de 4.2 mm (EETT: 4 mm) & Revisor\\
3.6 & Moldaje cadena & \num{22.48} & \num{23.77} & Faltaban fondos bajo dinteles y tapas de extremos & Revisor\\
8.1 & Bajadas AALL & \num{12.60} & \num{15.60} & Se tomó la cota del alero como largo de bajada & Revisor\\
3.1 & Ladrillo gran titán & \num{78.07} & \num{88.56} & Vanos descontados al 100\,\%; NCh 353 no descuenta los menores a \SI{1.5}{m^2} & NCh 353, revisor\\
3.5 & Enfierradura cadena & \num{173.65} & \num{194.77} & Sin traslapos de 40\,$d$ ni +5\,\% de patas & NCh 353, revisor\\
\bottomrule
\end{tabularx}
\caption{Correcciones a la cubicación propia tras la revisión en frío y la lectura de NCh 353.}
\label{tab:corr}
\end{table}
